\section{Conclusion}
\label{sec:conclusion}

UVIF proposes an explicit decomposition of five regulatory demands relevant to persistent cognitive systems: information acquisition, selective protection, computation cost, effective action, and future operating capacity. The present formulation distinguishes these soft influences from hard admissibility requirements and separates stationary equilibrium from invariant operation and adaptation under changing conditions.

The analytical example demonstrates stability in one restricted model and also shows that its dynamics can minimize a scalar objective. The contribution is therefore a proposed organization of modelling responsibilities, not a demonstrated replacement for optimization. Its relationship to active inference, resource-rational analysis, allostasis, and cognitive architectures is potentially complementary.

The decisive next step is a controlled comparison in which the components are independently specified, their computational costs are counted, and their predictions are tested against capable alternatives. Until such evidence is available, UVIF should be treated as a conceptual research programme whose usefulness, component structure, and range of applicability remain open to revision.
